% Responsable: P4
% Tablas y figuras: python -m evaluation.metrics  (results/tablas/*.tex, results/figures/*.png).
% \tablagenerada está definida en 05_set_de_prueba.tex.
\section{Resultados}
\label{sec:resultados}

\newcommand{\figuragenerada}[2][\linewidth]{%
  \IfFileExists{../../results/figures/#2.png}{\includegraphics[width=#1]{#2}}{%
    \pendiente{generar results/figures/#2.png}}}

\pendiente{P4: estas tablas solo tienen el parser propio. Falta volver a correr spaCy y
Transformers (P1) con el conjunto de 73 oraciones y el LLM (P5); después se regeneran con
python -m evaluation.metrics y se completa la discusión.}

\subsection{Aceptación}

\begin{table}[h]
  \centering
  \small
  \tablagenerada{aceptacion}\par\medskip
  \tablagenerada{aceptacion_espanol}
  \caption{Aceptación contra \texttt{en\_gramatica} (arriba) y contra \texttt{espanol\_valido}
    (abajo). ``---'': la métrica no aplica al sistema.}
  \label{tab:aceptacion}
\end{table}

\begin{figure}[h]
  \centering
  \figuragenerada{aceptacion}
  \caption{Métricas de aceptación contra las reglas de la gramática.}
  \label{fig:aceptacion}
\end{figure}

El parser propio acepta exactamente las 45 oraciones de la gramática y rechaza las 28 que no
lo están. Además, todas las rechaza por la razón esperada: 15 por error sintáctico y 13 por
concordancia (tabla~\ref{tab:tipo_error}). Este resultado no mide qué tan bueno es el enfoque,
porque el conjunto se etiquetó con la misma gramática que implementa el parser. Lo que muestra
es que la implementación es fiel a la especificación, incluso en los casos difíciles: la
concordancia a distancia (\emph{El perro de los niños come} se acepta y \emph{*...comen} se
rechaza) y la palabra ambigua \emph{bajo}.

La comparación interesante es contra \texttt{espanol\_valido}. Ahí el parser tiene precisión
1.00, pero su cobertura baja a 0.80. Nunca acepta español incorrecto, pero rechaza 11 de las 56
oraciones correctas: todo lo que la CFG no genera (sujeto omitido, orden libre, coordinación,
subordinadas y palabras fuera del léxico). Ese 20\,\% es la medida concreta de cuánto español
correcto queda fuera de la gramática.

\begin{figure}[h]
  \centering
  \figuragenerada[0.7\linewidth]{reglas_vs_espanol}
  \caption{Exactitud de la aceptación contra las dos referencias.}
  \label{fig:reglas_vs_espanol}
\end{figure}

\pendiente{P4 con P5: comparar con Transformers y el LLM. Si su exactitud contra
espanol\_valido supera a la de en\_gramatica, siguen su conocimiento del español y no las reglas.}

\begin{table}[h]
  \centering
  \small
  \resizebox{\linewidth}{!}{\tablagenerada{tipo_error}}
  \caption{Tipo de error esperado frente a la respuesta del sistema, en las 28 oraciones fuera
    de la gramática.}
  \label{tab:tipo_error}
\end{table}

\subsection{Resultados por nivel}

\begin{table}[h]
  \centering
  \small
  \tablagenerada{exactitud_aceptacion_por_nivel}\par\medskip
  \tablagenerada{etiquetado_por_nivel}
  \caption{Exactitud de la aceptación (arriba) y del etiquetado (abajo) por nivel.}
  \label{tab:por_nivel}
\end{table}

\begin{figure}[h]
  \centering
  \figuragenerada{por_nivel}
  \caption{Exactitud por nivel de complejidad.}
  \label{fig:por_nivel}
\end{figure}

\pendiente{P4: describir cómo cambia cada sistema del nivel 1 al 3.}

\subsection{Etiquetado y funciones sintácticas}

\begin{table}[h]
  \centering
  \small
  \tablagenerada{etiquetado}\quad\tablagenerada{funciones}
  \caption{Exactitud del etiquetado (y oraciones con otro número de palabras), y de las
    funciones sintácticas.}
  \label{tab:etiquetado}
\end{table}

\begin{table}[h]
  \centering
  \small
  \tablagenerada{etiquetado_categoria}
  \caption{Cobertura (\emph{recall}) del etiquetado por categoría.}
  \label{tab:etiquetado_categoria}
\end{table}

\begin{figure}[h]
  \centering
  \figuragenerada[0.8\linewidth]{etiquetado_categoria}
  \caption{Cobertura del etiquetado por categoría.}
  \label{fig:etiquetado_categoria}
\end{figure}

El parser etiqueta sin errores, y da bien el sujeto, el verbo y el objeto, en todas las
oraciones que puede derivar (58 de 73: las aceptadas y las que fallan solo en la
concordancia). En las otras 15 no produce etiquetas. Un etiquetador estadístico sí puede
analizarlas, y esa diferencia de cobertura es la que se discute en la
sección~\ref{sec:comparacion_nlp}.

\subsection{Unión de la frase preposicional}

\begin{table}[h]
  \centering
  \small
  \tablagenerada{pp_adjunta}\par\medskip
  \tablagenerada{pp_confusion}
  \caption{Unión de la primera frase preposicional en las 17 oraciones que la tienen.}
  \label{tab:pp}
\end{table}

\begin{figure}[h]
  \centering
  \figuragenerada{funciones_pp}
  \caption{Funciones sintácticas y unión de la frase preposicional.}
  \label{fig:funciones_pp}
\end{figure}

Aquí aparece el costo del conflicto LL(1) de la sección~\ref{sec:gramatica}. El parser une la
frase preposicional al sustantivo más cercano siempre que hay un objeto. Por eso acierta en
las 6 oraciones donde la unión correcta es al sustantivo, y falla en las 4 donde un hablante la
une al verbo (\emph{Ellos ven los árboles altos desde la casa}). Las 2 oraciones ambiguas las
resuelve sin avisar, como si no lo fueran, así que detecta 0 de 2 ambigüedades. Su exactitud es
0.59 (10 de 17). Una de las fallas es \emph{El perro corre en el parque}, que no se puede
derivar porque \emph{corre} no está en el léxico.

\subsection{Costo}

\begin{table}[h]
  \centering
  \small
  \tablagenerada{costo}
  \caption{Tiempo medio por oración y fallos de ejecución.}
  \label{tab:costo}
\end{table}

\pendiente{P4: comparar el tiempo del parser (centésimas de milisegundo por oración) con el de
spaCy, Transformers y el LLM cuando estén sus resultados.}
